A central methodological problem in modern LLM evaluation is \emph{attribution}: when a model answers correctly, did it generalize or did it retrieve a near-copy of the answer from its pretraining data? The recently released \talkienineteen~\citep{talkielm2026} is a 13B-parameter LLM trained on 260B tokens of strictly pre-1931 English text, paired with an architecture-matched modern twin \talkieweb trained on FineWeb. We test the meta-hypothesis that this hard 1930 corpus cutoff makes generalization tests \emph{qualitatively} easier --- not by reducing required sample sizes, but by making each correct answer on a post-1930 task uncontaminated by construction. We run three matched probes on \{\talkienineteen, \talkieweb, \gptfour\}: a 50-item post-1930 probe battery, a date-stratified \mmlu subsample, and \humaneval. On the probe battery, \talkienineteen attains 3\% post-1930 fact recall versus 73\% and 80\% for the modern baselines, with a paired per-token surprisal gap of +6.84 \bpt over its modern twin (paired permutation $p < 0.001$, $n=30$). Yet \talkienineteen successfully extends Python from a single in-context worked example on 4 of 11 items despite zero prior Python exposure --- uncontaminated generalization in the wild. Under a worst-case Bayesian attribution analysis, each correct \talkienineteen answer on a post-1930 item is worth $\approx 34$ bits of evidence for generalization, while the modern baselines yield $\le 0$ bits without an additional contamination audit. The advantage is qualitative: \talkienineteen turns generalization claims into unambiguous claims that no modern-LLM evaluation can produce without expensive corpus accounting.
